\solchapter{gr}

\sol{gr}{11} Linearize: $u^2\approx(GM/h^2)^2(1+2e\cos\omega\phi)$. Matching the $\cos\omega\phi$ terms fixes $\omega$; the perihelion recurs when $\omega\phi = 2\pi$.

\sol{gr}{12} Compare the light ray $x = ct - D$ with the asymptote $x = ct - c^2/a$ of the hyperbola. For the second part, the acceleration needed to hover at $r$ is $GM/r^2\sqrt{1-r_s/r}$; redshift the local Unruh temperature out to infinity.

\solchapter{astro}

\sol{astro}{11} Write both pressures in terms of $M$ and $R$. If they scale the same way with $R$, their ratio depends only on $M$. What happens if gravity wins at one radius?

\sol{astro}{12} For the lower limit, add the electron degeneracy energy $\sim\hbar^2n^{2/3}/m_e$ to the virial balance and find the maximum temperature the ball ever reaches. For the upper limit, compare $L\propto M^{3.5}$ with $L_{\text{Edd}}\propto M$.

\sol{astro}{13} Angular momentum $\propto MR^2/P$ and flux $\propto BR^2$ are conserved. For the maximum spin, set the centripetal acceleration at the equator equal to gravity.

\solchapter{cosmo}

\sol{cosmo}{9} After freeze-out, the neutrons decay as $e^{-t/880\un{s}}$. Each helium-4 nucleus holds two neutrons, so all the neutrons, and as many protons, end up in helium.

\sol{cosmo}{10} Multiply the growth factors of the two eras: $(a_{\text{eq}}/a)^2$ and then $1/a_{\text{eq}}$, with $a_{\text{eq}}\approx1/3400$. Inflation must shrink $\abs{\Omega - 1}$ from about $1$ to this value.

\sol{cosmo}{11} Find the fraction of the sky covered by the stars in a thin shell at distance $r$. Does it depend on $r$? For the second part, compare the distance a line of sight must go to hit a star with the distance light has travelled since the Big Bang.

\solchapter{gw}

\sol{gw}{9} Differentiate $E = -Gm_1m_2/2a$ to get $\dot a$, then separate variables and integrate $a^3\,\dd a$.

\sol{gw}{10} Write $a$ in terms of $f$ using Kepler's law with $f_{\text{orb}} = f/2$, then $\dot f/f = -\tfrac32\dot a/a$. Collect the masses into $m_1m_2(m_1+m_2)^{-1/3}$.

\sol{gw}{11} $H_0 = v/d$. For the error, think about what else besides distance changes the amplitude of the wave, and how errors on independent events combine.
